% Baseline evidence: ch4/papers/RAG/main.tex, reviewed 2026-09-09.
\section{From sparse lexical experience to retrievable evidence}
\label{sec:ch4:thesis-introduction}

Learning a word involves more than making its form available for production. A learner must also use it in combinations that respect grammatical dependencies, including when the word has been encountered infrequently. Chapter~\ref{chp2_chap} examined the first problem through expressive vocabulary: under the specified exposure and generation budgets, low-frequency words were less reliably represented in model output. The present chapter turns to a related question about syntactic use. When a predictive model assigns unreliable probabilities to sentences containing infrequent content words, can access to particular previous instances improve its sensitivity to grammatical contrasts?

The link between these questions is the availability of evidence from sparse experience. It is not a direct continuation of the preceding experiment. The vocabulary study used character-level learners and a production-based criterion, whereas this study uses GPT-2 models and probability comparisons between grammatical and ungrammatical sentences. Their corpora and training histories also differ. Retrieval is therefore tested here as a possible resource for frequency-sensitive linguistic performance, without assuming that it resolves the expressive-vocabulary discrepancy established in the previous chapter.

The distinction between learning from experience and retaining access to experience motivates the intervention. Parametric learning changes shared weights across many examples; retrieval allows a particular stored context to contribute directly to a later prediction. Complementary Learning Systems theory provides a cognitive motivation for investigating such a division of labour between gradually acquired regularities and access to specific experiences \citep{mcclelland1995there,kumaran2016learning}. The hypothesis examined here is functional: when the parametric prediction is unreliable, a relevant stored instance may supply useful evidence. The computational implementation does not reproduce hippocampal encoding or establish the neural basis of human grammatical processing. It makes one proposed memory function experimentally manipulable.

The study implements retrieval augmentation through a $k$-nearest-neighbour language model ($k$NN-LM; \citealt{khandelwal2020knnlm}). A datastore associates representations of previous contexts with their following tokens. At inference time, nearest-neighbour search supplies a token distribution that is interpolated with the base model's distribution, with retrieval weight $\lambda=0.25$ in the reported experiments. This form of retrieval-augmented language modelling changes the evidence used to score a sentence. Its immediate effect is measured without retraining the base model, so improved performance can reflect better access to stored evidence without demonstrating that new grammatical knowledge has been consolidated into the model's parameters.

Two settings test the intervention: GPT-2 Small trained from scratch on a 50.03M-word mixture of child-directed speech, conversational speech, stories, and subtitles, and the publicly pretrained GPT-2 XL model. The smaller setting constrains the base learner's training corpus, while the larger setting asks whether a similar intervention remains useful under broad pretraining. They are separate tests within different systems, rather than a controlled manipulation of data scale alone. The datastore and any externally trained representations must also be distinguished from the base learner's training input when interpreting the amount and source of available experience.
% AUTHOR CHECK: Document datastore corpus/split/size and its overlap with training and evaluation; current methods specify GPT-2 keys, not DPR.

The experiments compare each base model with its retrieval-augmented counterpart on frequency-stratified subject--verb agreement, wh-question, and relative-clause contrasts selected from BLiMP, Zorro, and BIG-bench. The central comparison is whether retrieval improves low-frequency items more than high-frequency items within the same model and corpus. Further analyses vary the representations used for retrieval, the presence of examples containing the target syntactic phenomenon, and retrieval granularity, context size, and neighbour count. These comparisons ask both whether access helps and what makes a stored instance useful.

The reported results show selective but incomplete compensation. Retrieval narrows the high--low accuracy gap in all three phenomenon groups and both model settings, with particularly large absolute gains for relative clauses. The gain depends on representation and retrieval configuration; simply making more neighbours available does not produce uniformly better performance. The article below presents these results and their technical analyses. The chapter-level discussion then considers their contribution to the thesis: a controlled demonstration that access to stored instances can improve specific linguistic decisions, together with a more precise account of the evidence needed to connect that effect to language learning and memory.
